\section{A prime-field linear rank cutoff with the whole partner retained}
\label{sec:prime-field-linear-point-cutoff}

\paragraph{Author freeze and provenance.}
D-TD65, question Q2fu, has complete author proof and source
readings and a first exact certificate PASS. Complete root and
independent whole-proof reviews also pass.
The author combines the accepted degree barrier
D-TD58, the complete projective-factor interfaces, and the supplied
published character and genus proofs specified below. The quantitative
genus extraction and graph-coset count are written explicitly.
No independent-origin claim is made for the accepted source,
normalizer, tensor, or full-period Fourier arguments.

Put
\[
 D_0=10^{64},\qquad L_0=1100000,\qquad T_0=2^{L_0},
 \qquad B=10^{30},\qquad \epsilon=2^{-B},
 \qquad C_0=2\cdot10^{30},\qquad J_0=2C_0+1.
 \tag{PFL.1}
\]
A projective support is the degree minus the largest eigenspace
dimension, with all eigenvalue targets allowed.

\begin{theorem}[D-TD65]
Let \(C\to\mathbf P^1\) be an actual connected finite \(J\)-Galois
cover with a faithful irreducible complex coefficient module \(V\)
of degree \(e\), and a generating product-one inertia tuple satisfying
\[
                    \sum_i\operatorname{rank}(1-T_i)=2e+6.
 \tag{PFL.2}
\]
Suppose an actual stable whole factorization
\(V=W\otimes U\) has the following properties. Its finite projective
image \(P\) on \(W\) contains and normalizes the projective image of
a finite perfect central product \(E_1\cdots E_t\), acting faithfully
and irreducibly on \(W\). Its commuting quasisimple components are
permuted transitively by \(P\). Each \(E_j\) is an actual central
quotient of \(\operatorname{SL}_d(p)\), where \(p\) is prime and
\(d\ge D_0\). Write \(W=\bigotimes_{j=1}^t W_j\), where the faithful
irreducible component degrees have their common value \(a\).
If \(1\le t\le T_0\), this configuration is impossible.

The entire external partner \(U\), including its dimension, image
and scalar phases, is unrestricted. No number of projected inertia
points, rationality of a component-stabilizer base, character field,
or self-duality hypothesis is imposed.

Consequently, in the minimal naturally nonsplit exact-imaginary-
quadratic source class of D-TD54, every layer of any fixed complete
D-TD12 recursion having the indicated prime-field linear
named-cover presentation has \(d<D_0\).
\end{theorem}

\subsection*{1. A quantitative consequence of the published genus proof}

We first record precisely the stronger consequence extracted from
\cite{PFL:FGM2014}. If
\(\operatorname{PSL}_d(p)\le G\le\operatorname{PGL}_d(p)\),
\(p\) is prime and \(d\ge32\), then every actual generating
product-one tuple on the projective points, of degree
\(N=(p^d-1)/(p-1)\), has quotient genus \(\gamma\) satisfying
\[
                         \gamma\ge1+\frac{N}{10000}.
 \tag{PFL.3}
\]
This is a consequence of the supplied proof, not the literal
statement of that paper's Theorem 5.

Remove identity tuple entries. The actual regular \(G\)-cover of
the base sphere has genus at least two by D-TD02: its group contains
the nonsolvable \(\operatorname{PSL}_d(p)\) and has order greater
than 60, excluding the genus-zero finite groups and the solvable
elliptic finite groups. Its Riemann--Hurwitz equality therefore gives
\(A(\mathbf o)=\sum_i(1-1/o(x_i))>2\), and in particular \(r\ge3\).
This supplies the positive-signature-excess hypothesis used by the
published calculations to exclude spherical and Euclidean signatures.

Here is the full parameter and system check. For a projective
operator \(x\), let \(v(x)\) be the codimension of the largest
\(\mathbf F_p\)-eigenspace among its scalar lifts, using dimension
zero when no rational eigenvalue occurs. This is the source's
verbal definition and the convention used in its Fact 8;
we do not use the conflicting displayed maximum-codimension
expression on printed p.242. Put
\(\kappa(x)=o(x)^{-1}\sum_{j=0}^{o(x)-1}p^{-v(x^j)}\).
The complete proof of Proposition 36, printed pp.262--263, gives
\(f_\Omega(x)<p^{-v(x)}+p^{-d/2}\) for every nonidentity
projective operator. In particular the Grassmann error is less
than \(1/100\).

For clarity, the system inequalities used in Section 2 do not
require a quotient to have the exact index asserted in Lemma 15's
wording. Take a spherical signature \(\mathbf e\) from that lemma.
Its universal finite spherical group has order
\(L=2/(2-A(\mathbf e))\), where
\(A(\mathbf e)=\sum_i(1-1/e_i)\).
The regular sphere cover has \(L/e_i\)
points over its \(i\)-th mark. The punctured cover therefore has a
product-one generating set consisting of those conjugates of the
\(e_i\)-th powers of the original loops. Map this construction into
the actual group generated by the given tuple. The resulting group
\(K\) is normal and the actual quotient is a \emph{quotient} of the
cyclic, dihedral, \(A_4,S_4\), or \(A_5\) spherical group; it need
not have order \(L\). The conjugate multiplicities \(L/e_i\) and
the product-one relation are unchanged, with repetitions allowed.
This is exactly the information needed below. Existence and the
sphere genus here follow either from the usual explicit spherical
actions or from their regular permutation tuples and the accepted
Riemann-existence/Riemann--Hurwitz interface of D-TD02.

The huge simple normal subgroup \(\operatorname{PSL}_d(p)\) maps
trivially to each such quotient: the cyclic and dihedral groups
are solvable, and the other three have order at most 60. Thus
\(K\) contains that subgroup. Any scalar-chosen linear lift of its
generators acts irreducibly on the natural \(d\)-space, with no
invariant vectors or covectors.

The Scott inequality in the source has a short direct verification.
For linear \(g_1\cdots g_s=1\), put \(V_i=\operatorname{im}(g_i-1)\)
and map \(\bigoplus_iV_i\) to \(V\) by
\((w_i)\mapsto\sum_i g_1\cdots g_{i-1}w_i\).
Its image has codimension \(\dim(V^*)^{\langle g_i\rangle}\).
Its kernel contains the image of
\(v\mapsto((g_i-1)v)_i\), of dimension
\(d-\dim V^{\langle g_i\rangle}\). Hence
\(\sum_i\operatorname{rank}(g_i-1)\ge
2d-\dim V^{\langle g_i\rangle}-\dim(V^*)^{\langle g_i\rangle}\).
Choose all but one of the spherical lifted generators to attain
their largest rational eigenspace and choose the last to make the
product one. Its rank is at most \(d\). This gives precisely
Lemma 16's inequalities, including its stronger \(2d\) inequality
at \(p=2\), where the scalar group is trivial. All subsequent
system estimates in Lemmas 17--24 therefore apply with their
stated constants.

Suppose now that (PFL.3) fails. The source's normalized genus
parameter obeys
\[
 \epsilon_0=\frac{2(\gamma-1)}{N}<0.0002,\qquad
 \sum_i\kappa(x_i)>
 r-2-0.01\sum_i(1-1/o(x_i))-0.0002.
 \tag{PFL.4}
\]
These are exactly the numerical hypotheses used in Fact 25 and
the large-dimension calculations of Sections 2.3--2.4,
printed pp.249--261. The numerical conclusions of Lemma 26
are already satisfied by \(d\ge32\). Its separate \(d=2\)
genus argument, and Lemma 29's separate \(d=3\) genus argument,
are not used. Every remaining calculation uses (PFL.4), the
individual-power estimates of Section 2.1 and the system
inequalities just checked. Lemma 30 reduces to \(p\le19\);
Propositions 31--35 give exactly Table 1, whose largest natural
prime-field dimension is 21. This contradicts \(d\ge32\) and
proves (PFL.3). The supplied extensive numerical verification is
a published input; no independent rerun of its GAP computations
is claimed. The prime-field restriction is explicit.

We will need (PFL.3) over arbitrary actual intermediate bases.
For a connected point cover \(Y\to X\) of degree \(N\), with
monodromy containing \(\operatorname{PSL}_d(p)\), it implies
\[
                            g(Y)\ge1+\frac{N}{10000}
                            \quad\text{for every }g(X)\ge0.
 \tag{PFL.5}
\]
For \(g(X)=0\) this is (PFL.3). For \(g(X)=1\), an unramified
finite monodromy image would be abelian, so there is a nontrivial
branch permutation. Its fixed-point ratio is at most
\((p^{d-1}+p-2)/(p^d-1)\le2/3\): decompose its rational
eigenspaces and maximize the number of lines by dimensions
\(d-1\) and 1. Its permutation index is at least \(N/6\).
Riemann--Hurwitz gives \(g(Y)-1\ge N/12\).
For \(g(X)\ge2\), it gives \(g(Y)-1\ge N\).
These are both stronger than (PFL.5).

\subsection*{2. The actual full prime-field normalizer}

The actual projective normalizer \(P\) embeds by conjugation into
\(\operatorname{Aut}(\operatorname{PSL}_d(p)^t)\).
Indeed an operator acting trivially on this central quotient has
central commutators with the represented perfect core. The
central discrepancy is a homomorphism and hence trivial; Schur's
lemma then makes the operator scalar. This is the exact accepted
D-TD58 perfectness argument, applied before taking projective image.

Steinberg's Theorem 3.2 and its supplied normal-type proof,
\cite[pp.607--613, Sections 3--6]{PFL:Steinberg1960}, say that
every automorphism of \(\operatorname{PSL}_d(p)\) is inner,
diagonal, field and graph. Here the field part is trivial.
Inner-diagonal automorphisms are represented by
\(\operatorname{PGL}_d(p)\), and the sole nontrivial graph coset
is inverse-transpose times an inner-diagonal automorphism.
The inverse-transpose realization follows directly on the
elementary root matrices. This provides an actual action, for
each component, on the disjoint union of its points and dual
points. Graph automorphisms interchange the two types; linear
ones preserve them. Transporting the named-cover presentation
through the actual component isomorphisms supplies these actions
on all components, without assuming a separate quotient map
to any permuted component.

Thus \(J\) has an actual permutation action on
\[
                  \Omega=\coprod_{j=1}^t
                   \bigl(\mathbf P^{d-1}(\mathbf F_p)
                          \sqcup\mathbf P^{d-1}(\mathbf F_p)^*\bigr),
                  \qquad |\Omega|=2tN.
 \tag{PFL.6}
\]
There are one or two orbits. For each orbit, the stabilizer of
an oriented component gives an actual intermediate curve \(X\).
Its local image is a subgroup of \(\operatorname{PGL}_d(p)\)
containing \(\operatorname{PSL}_d(p)\). The corresponding point
cover \(Y\to X\) is connected of degree \(N\). It is the same
curve as that orbit's quotient cover over the original sphere.
Indeed both are the quotient by the inverse image of a point
stabilizer. Apply (PFL.5) to each such curve, and sum the actual
Riemann--Hurwitz equalities. If \(y_i\) is the image of the
original \(T_i\), then
\[
               \sum_i\frac{\operatorname{ind}_\Omega(y_i)}{2tN}
                          \ge2+\frac1{10000t}.
 \tag{PFL.7}
\]
No rationality hypothesis has been placed on those intermediate
component or orientation bases.

\subsection*{3. Uniform inner-diagonal traces at bounded level}

The transitive version of D-TD58 gives \(a<p^{C_0d}\).
Inflate an irreducible component character to
\(\operatorname{SL}_d(p)\), and let \(j\) be its level.
The exact GLT degree and extension inputs are the accepted
D-LNW-01, together with
\cite[Definition 2, Lemma 7.2, Theorem 7.3 and Corollary 8.6,
pp.51--53,67--68]{UNW:GLT2020}.
They give
\[
                1\le j\le J_0<d/2,\qquad a>p^{d/12},
 \tag{PFL.8}
\]
and an irreducible general-linear extension of the same degree.
For completeness, \(j\ge d/2\) would give
\(a>(2/3)p^{d^2/4-3}>p^{C_0d}\).
For \(j<d/2\), \(a\ge p^{j(d-j)}/(p-1)>
p^{jd/2-1}\), so \(j\ge2C_0+2\) is impossible.
Also \(j(d-j)\ge d-1\), giving \(a>p^{d-2}>p^{d/12}\).
True level is selected by a determinant twist, which does not
affect absolute character values.

For \(j\ge2\), use the stronger displayed estimate (8.11)
and its complete dual-pair derivation on published pp.57--61,
rather than only the numerical corollary. If \(s\) is geometric
natural support and \(F=\alpha(1)p^{dj}/|\operatorname{GL}_j(p)|\),
it gives
\[
 \begin{split}
 a&\ge F(1-E),&
 |\widetilde\chi(g)|&\le F(p^{-js}+E),\\
 E&=|\operatorname{GL}_j(p)|
       (1+\sqrt{8p^{j^2+j}})p^{-d}
       \le4p^{2j^2-d}\le p^{-d/2}.
 \end{split}
 \tag{PFL.9}
\]
The middle term of (8.11) is absorbed because
\((d-1)(j-1)+1\le d(j-1)\).
All constants follow from \(d\ge D_0\), \(j\le J_0\).
Consequently the normalized ratio is at most \(p^{-s}+4p^{-d/2}\).

If a largest geometric eigenspace is rational, its projective
points show \(p^{-s}\le f_{\rm pt}(g)+p^{-d}\).
If it is not rational, its Galois orbit has at least two members,
so \(s\ge d/2\). Therefore the normalized ratio is at most
\(f_{\rm pt}(g)+5p^{-d/2}\).

Here is also a direct treatment of level one. By its definition,
a determinant twist of the character occurs in the natural-vector
permutation module. Write the nonzero-vector action as induction
from a vector stabilizer. Induction through its line stabilizer
decomposes it over the \(p-1\) scalar characters \(\theta\) into
monomial line-bundle modules of degree \(N\). A nontrivial
\(\theta\) induces irreducibly: the line-stabilizer double-coset
decomposition has two terms, and on the off-diagonal intersection
the two line scalars vary independently, so its intertwining
space is zero. For \(\theta=1\), double transitivity gives
\(1+\psi\) with \(\psi\) irreducible of degree \(N-1\).
These exhaust the nontrivial constituents of true level one.
The nontrivial line bundles have absolute trace at most the
number of fixed lines. For \(\psi\), the ratio is at most
\(f_{\rm pt}(g)\) when that number is positive, and at most
\(1/(N-1)\le p^{-d/2}\) when it is zero.
Restriction remains irreducible by the cited Corollary 8.6.

If an actual component intertwiner induces the inner-diagonal
automorphism represented by \(g\), it is a scalar multiple of
\(\widetilde\rho(g)\). Perfectness removes any central discrepancy
and Schur's lemma gives the scalar. With unitary normalizations
the scalar has modulus one. We have proved, for every such
actual operator \(A\),
\[
           \frac{|\operatorname{tr}A|}{a}
                \le f_{\rm pt}(g)+5p^{-d/2}
                \le f_{\rm pt}(g)+\epsilon.
 \tag{PFL.10}
\]
Point and dual-point fixed fractions agree for a linear operator,
since its transpose has the same eigenspace dimensions.

\subsection*{4. Every compatible graph-coset operator}

Let \(A\) induce a graph-diagonal automorphism
\(\alpha(X)=M X^{-T}M^{-1}\) on
\(S=\operatorname{PSL}_d(p)\). Its projectively fixed lifts in
\(\operatorname{SL}_d(p)\) satisfy
\(X^T B_1X=cB_1\) for a nonsingular \(B_1=M^{-1}\) and a scalar
\(c\in\mathbf F_p^*\). No symmetry of \(B_1\) is assumed.
For independent columns \(x_1,\ldots,x_d\), the first \(j-1\)
conditions \(x_i^TB_1x_j=c(B_1)_{ij}\) impose independent linear
equations on column \(j\). Ignoring all other equations gives at
most \(p^{d(d+1)/2}\) matrices for each \(c\).
Since
\(|\operatorname{SL}_d(p)|=
p^{d^2-1}\prod_{i=2}^d(1-p^{-i})>p^{d^2-1}/2\), it follows that
\[
                   \frac{|C_S(\alpha)|}{|S|}
                     \le2p^{-d^2/2+d/2+2}.
 \tag{PFL.11}
\]
This counts the entire projective fixed subgroup, including
possible nontrivial scalar discrepancies.

Put \(R=\lfloor d/8\rfloor\). A matrix of geometric support
less than \(R\) has a unique largest rational eigenvalue:
its dimension is greater than \(d/2\), so Galois cannot move it.
It has the form \(\lambda I+BC\), with
\(B\) a \(d\)-by-\(R\) matrix and \(C\) an \(R\)-by-\(d\) matrix.
Thus there are at most \(p^{2dR+1}\) possible projective elements
of this support, allowing harmless multiple counting.
Every nonempty fiber of \(x\mapsto\alpha(x)x^{-1}\) has
\(|C_S(\alpha)|\) elements. The proportion of low-support
commutators is therefore at most \(2p^{-d^2/8}\).

Use the all-element Larsen--Tiep inequality in (CFD.6), with
\(\sigma_0=10^{-29}\), on the inflated SL character.
Equation (PFL.8) bounds the other commutator ratios by
\(p^{-\sigma_0R/12}\). Schur averaging, exactly as in (CFD.10),
therefore gives
\[
 \frac{|\operatorname{tr}A|^2}{a^2}
   \le2p^{-d^2/8}+p^{-\sigma_0R/12}
   \le3p^{-2B-4}<p^{-2B}.
 \tag{PFL.12}
\]
Indeed \(R\ge d/9\), \(d^2/8\ge2B+4\), and
\(d/(108\cdot10^{29})\ge2B+4\).
These are deliberately weak exact inequalities at \(d\ge10^{64}\).
The trace is thus less than \(\epsilon\).
A graph operator fixes no point of the points-plus-dual-points
action, so (PFL.10) with that action's fixed fraction remains
valid for every full compatible component normalizer.

\subsection*{5. Tensor permutations and all original inertia periods}

For any representative \(A\) of the whole layer operator, write
\(\pi\) for its component permutation, \(c\) for the number of
fixed components, and \(I=t-\#\{\text{cycles of }\pi\}\).
The exact tensor-cycle trace formula bounds its normalized trace
by \(a^{-I}\) times the product of the normalized traces of the
fixed components. Return operators on the other cycles are
bounded by one after normalization.

If a fixed component has graph type, this product is at most
\(\epsilon\). Otherwise all fixed components have the bound
(PFL.10). Their true normalized ratios are in \([0,1]\), and
their product is at most their average. When \(c\ge1\),
\(a^I\ge1+2I\ge t-c+1\ge t/c\).
Thus the whole trace is at most \(f_\Omega(y)+\epsilon\).
If \(c=0\), then \(I\ge1\) and the trace is at most
\(1/a<\epsilon\), while \(f_\Omega(y)=0\).
The conclusion holds also for \(t=1\). A projectively trivial
whole-layer operator is scalar, and its ratio and fixed fraction
are both exactly one.

Only at most 160 original inertia entries require an error term.
The accepted full-normalizer floor in HD-FT20--21 gives
\(\nu(A)/\dim W\ge1/80\) on every nonscalar layer operator.
It applies to this whole central product and includes its
permuting normalizers. More explicitly, a fixed nonscalar component
has the HD-FT20 floor \(1/80\); a moved component cycle has the
HD-FT21 quarter floor. Spectral convolution gives the asserted
whole-layer floor in either case.
Tensor convolution preserves the floor with the entire \(U\).
Since \(e\ge a>480\), (PFL.2) implies that the number \(k\) of
original entries with nonidentity projective layer image satisfies
\[
                  k/80\le2+6/e<161/80,\qquad k\le160.
 \tag{PFL.13}
\]

Let \(L_i\) be the full original order of \(T_i\). Reciprocal
scalar normalization of its two tensor factors makes both
\(L_i\)-th powers the identity, without changing \(T_i\).
Every power has the just-proved trace bound, and the actual
permutation \(y_i\) has order dividing \(L_i\).
The full \(L_i\)-power Fourier average at eigenvalue one gives
\[
 \frac{\operatorname{rank}(1-T_i)}e
       \ge\frac{\operatorname{ind}_\Omega(y_i)}{|\Omega|}
                 -\epsilon\quad(y_i\ne1).
 \tag{PFL.14}
\]
If \(y_i=1\), retain the nonnegative actual rank without an error.
There is no replacement of an original period by a projected
one and no restriction on the other factor's phases.

Sum (PFL.14), use (PFL.7) and (PFL.13), and obtain
\[
           2+\frac6e\ge2+\frac1{10000t}-160\epsilon.
 \tag{PFL.15}
\]
But \(e\ge a>2^{d/12}\), so \(6/e<6\epsilon\). Also
\(t\le2^{L_0}\) and \(B>L_0+23\) imply
\[
             166\epsilon<2^{-L_0-14}
                            <\frac1{10000t}.
 \tag{PFL.16}
\]
This contradicts (PFL.15), proving the standalone assertion.

\subsection*{6. The exact minimal-source consequence}

Fix the same complete primitive recursion as in D-TD54.
Its accepted cutoff \(N_*\) has the elementary upper estimate
\[
 \begin{gathered}
 240!<2^{1920},\qquad M_*<2^{2168},\qquad F_*<2^{2177},\\
 N_*=[100(1+M_*)F_*^{240}]^2
                <2^{1049312}<2^{1100000}=T_0.
 \end{gathered}
 \tag{PFL.17}
\]
The TD54 whole-source floor is \(1/M_0\), with
\(M_0=\max(80,N_*-2)<T_0\). Applying D-TD16 in its actual
original-source scope gives \(t\le\lfloor M_0/2\rfloor<T_0\)
for every transitive component orbit. D-TD12 supplies the
irreducible represented central product and its actual quotient
normalizer, while retaining all earlier and later pieces in
the same entire \(U\). The preceding theorem therefore excludes
\(d\ge D_0\) for every prime-field linear layer with the named
central-quotient presentation. No independent subgroups of the
original \(J\) are presumed for later extraction cores.

This is a genuine bound on an unbounded natural-rank residual
left by D-TD58's degree restriction. It does not extend the
quantitative FGM calculation to proper prime powers, treat
unitary or symplectic low-level characters, bound all Lie
parameters or original overgroups, or decide source existence.
The standalone theorem needs only (PFL.2) and its stated actual
factor assumptions; the corollary retains exactly D-TD54's
minimal, exact-field and naturally nonsplit hypotheses.

